\documentclass[a4paper]{article}
\usepackage[pdftex]{graphicx}
\usepackage[utf8]{inputenc}
\usepackage{enumerate}
\usepackage{icomma}
\usepackage{siunitx}
\sisetup{locale=DE}
\usepackage{amssymb}
\usepackage{tikz}
\usepackage{href-ul}
\hypersetup{
	colorlinks=true,
	linkcolor=blue,
	urlcolor=blue}
\usepackage{geometry}
\geometry{a4paper, top=15mm, left=15mm, right=15mm, bottom=15mm,
	headsep=10mm, footskip=12mm}

\begin{document}
	\begin{tikzpicture}(10,3)
		\draw[ultra thick](2,0) --(10,0) -- (10,1.3) --(2,1.3) -- (2,0);
		\draw[fill=black](2,0)-- (0,.6) -- (2,.6) -- (2,0);
		\node at (6,.6) {\large solution sheet from \href{http://www.okuyakl.com}{www.okuyakl.com}};
	\end{tikzpicture}
	\vspace{0.5 cm}
	
	\noindent
	\begin{minipage}{0.5\textwidth}
		\noindent{\bf Task 1. a)}\\
		\includegraphics[width=7 cm]{reins054}
	\end{minipage}
	\hfill
	\begin{minipage}{0.5\textwidth}
		\noindent{\bf Task 1. b)}\\
		$$\mathbb{G}=]0;12[$$
		\noindent{\bf Task 1. c)}\\
		$$U(x)=2 \cdot a + 2 \cdot b = 2 \cdot (12-x) + 2 \cdot (5+x)$$
		$$ = 24 - 2x + 10 + 2x = \SI{34}{\centi\meter}$$
		The scope always remains the same, it does not depend on $x$.
	\end{minipage}
	
	\noindent{\bf Task 1. d)}\\
	$$A(x)= a \cdot b = (12-x)\cdot (5+x) = -x^2 + 7x + 60$$
	\noindent{\bf Task 1. e)}\\
	Square complement gives:
	$$T_{max}=\SI{72,25}{\centi\meter^2} \quad \textnormal{for} \quad x=3.5$$
	
	\noindent
	\begin{minipage}{0.5\textwidth}
		\noindent{\bf Task 2. a)}\\
		\includegraphics[width=7 cm]{rzwei054}
	\end{minipage}
	\hfill
	\begin{minipage}{0.5\textwidth}
		\noindent{\bf Task 2. b)}\\
		$$\mathbb{G}=]0;6[$$
		\noindent{\bf Task 2. c)}\\
		$$U(x) = 2 \cdot (6-x) + 2 \cdot (8+2x)$$
		$$ = 12 - 2x + 16 + 4x = (28 + 2x)~\SI{}{\centi\meter}$$
	\end{minipage}
	
	\noindent{\bf Task 2. d)}\\
	$$A(x)= (6-x)\cdot (8+2x) = -2x^2 + 4x + 48$$
	\noindent{\bf Task 2. e)}\\
	Square complement gives:
	$$T_{max}=\SI{50}{\centi \meter^2} \quad \textnormal{for} \quad x=1$$
	
	\noindent
	\begin{minipage}{0.5\textwidth}
		\noindent{\bf Task 3. a) b)}\\
		\includegraphics[width=7 cm]{rdrei054}
	\end{minipage}
	\hfill
	\begin{minipage}{0.5\textwidth}
		Perimeter square: $U_Q = 4 \cdot a = 4 \cdot \SI{8}{\centi\meter} = \SI{32}{\centi\meter} $\\
		Area: $A_Q=a^2=(\SI{8}{\centi\meter})^2=\SI{64}{\centi\meter^2}$\\
		\noindent{\bf Task 3. c)}\\
		$$\mathbb{G}=]0;8[$$
		\noindent{\bf Task 3. d)}\\
		$$U(x) = 2 \cdot (8-x) + 2 \cdot (8+2x)$$
		$$ = (32 + 2x)~\SI{}{\centi\meter}$$
	\end{minipage}
	
	\noindent{\bf Task 3. e)}\\
	$$A(x)= (8-x) \cdot (8+2x) = -2x^2 + 8x + 64$$
	\noindent{\bf Task 3. f)}\\
	Square complement gives:
	$$T_{max}=\SI{72}{\centi \meter^2} \quad \textnormal{for} \quad x=2$$
	\vspace{0.5 cm}
	
	\noindent
	\begin{minipage}{0.5\textwidth}
		\noindent{\bf Task 4. a) b)}\\
		\includegraphics[width=7 cm]{dvier054}
	\end{minipage}
	\hfill
	\begin{minipage}{0.5\textwidth}
		Area: $A_D={1\over2}gh={1\over2}\SI{4}{\centi\meter}\cdot \SI{10}{\centi\meter}=\SI{20}{\centi \meter^2}$\\
		\noindent{\bf Task 4. c)}\\
		$$\mathbb{G}=]0;10[$$
		\noindent{\bf Task 4. d)}\\
		$$A(x)= {1\over2} (10-x) \cdot (4+2x) = -x^2 + 8x + 20$$
		\noindent{\bf Task 4. e)}\\
		Square complement gives:
		$$T_{max}=\SI{36}{\centi \meter^2} \quad \textnormal{for} \quad x=4$$
		
	\end{minipage}
	
	\vspace{0.5 cm}
	
	
	\noindent
	\begin{minipage}{0.5\textwidth}
		\noindent{\bf Task 5. a) b)}\\
		\includegraphics[width=7 cm]{rfunf054}
	\end{minipage}
	\hfill
	\begin{minipage}{0.5\textwidth}
		Area: $A_Q={1\over2}ef={1\over2}\SI{8}{\centi\meter}\cdot \SI{6}{\centi\meter}=\SI{24}{\centi \meter^2}$\\
		\noindent{\bf Task 5. c)}\\
		$$\mathbb{G}=]0;4[$$
		\noindent{\bf Task 5. d)}\\
		$$A(x)= {1\over2} (8-2x) \cdot (6+2x) = -2x^2 + 2x + 24$$
		\noindent{\bf Task 5. e)}\\
		Square complement gives:
		$$T_{max}=\SI{24,5}{\centi \meter^2} \quad \textnormal{for} \quad x={1\over2}$$
		This diamond is a square.
	\end{minipage}
	
	
	
	\vspace{0.5 cm}
	
	\begin{center}
		\includegraphics[width=7 cm]{../../viecher/eendcomic.pdf}
		
		Here you can go back to the \href{https://www.okuyakl.de/math/m8terL054/ae054.pdf}{task sheet}
	\end{center}
	
\end{document}